\section{Zukünftige Arbeiten}

Die vorliegende Arbeit hat ein umfassendes System zur Regelung eines selbstbalancierenden Fahrrads entwickelt, das eine duale Controller-Architektur mit PID-basierter Balance-Regelung und YOLO-gestützter Vision-Kontrolle implementiert. Trotz der erzielten Erfolge in der Simulation zeigen sich mehrere vielversprechende Forschungsrichtungen, die das System erheblich erweitern und verbessern könnten.

\subsection{Sensor-Fusion für robuste Zustandsschätzung}

Das aktuelle System basiert primär auf IMU-Daten zur Roll-Winkel-Bestimmung und Kamerabildern für die Spurerkennung. Eine Integration multipler Sensoren durch fortgeschrittene Sensor-Fusion-Techniken würde die Robustheit und Genauigkeit des Systems erheblich steigern.

\subsubsection{Erweiterte Kalman-Filter-Implementierung}
Die Implementierung eines Extended Kalman Filters (EKF) oder Unscented Kalman Filters (UKF) zur Fusion von IMU-, Kamera- und zusätzlichen Sensordaten stellt eine natürliche Erweiterung dar. Der EKF könnte folgende Sensormodalitäten integrieren:

\begin{itemize}
    \item \textbf{IMU-Fusion}: Kombination von Gyroskop-, Beschleunigungs- und Magnetometer-Daten für präzise Orientierungsschätzung
    \item \textbf{Optischer Fluss}: Ergänzung zur YOLO-Segmentierung für kontinuierliche Bewegungsschätzung
    \item \textbf{Ultraschall/LiDAR}: Abstandsmessung für Hinderniserkennung und 3D-Umgebungsmodellierung
    \item \textbf{Encoder-Fusion}: Integration von Rad- und Lenkwinkel-Encodern für kinematische Konsistenz
\end{itemize}

Die Zustandsschätzung würde dabei den Systemzustand $\mathbf{x} = [x, y, \theta, \dot{x}, \dot{y}, \dot{\theta}, \phi_{roll}, \delta_{steer}]^T$ umfassen, wobei die Prozess- und Messmodelle die nichtlineare Fahrradkinematik berücksichtigen.

\subsubsection{Adaptive Sensor-Gewichtung}
Ein adaptives Gewichtungssystem könnte die Zuverlässigkeit verschiedener Sensoren dynamisch bewerten und entsprechend anpassen. Bei schlechten Lichtverhältnissen würde beispielsweise die Gewichtung der Kameradaten reduziert und die IMU-Daten stärker berücksichtigt.

\subsection{Reinforcement Learning für optimale Regelung}

Während die aktuelle PID-Regelung für grundlegende Stabilisierung ausreicht, bietet Reinforcement Learning (RL) das Potenzial für adaptive und optimale Regelstrategien, die sich an verschiedene Fahrbedingungen anpassen können.

\subsubsection{Deep Deterministic Policy Gradient (DDPG)}
Die Implementierung eines DDPG-Algorithmus würde kontinuierliche Aktionsräume für Lenkwinkel und Geschwindigkeit ermöglichen. Der Zustandsraum könnte erweitert werden um:

\begin{itemize}
    \item Aktuelle und historische Sensordaten (IMU, Kamera, Geschwindigkeit)
    \item Pfadkrümmung und Vorausschau-Informationen
    \item Umgebungsbedingungen (Windstärke, Oberflächenbeschaffenheit)
    \item Systemzustand (Batteriestand, Motortemperatur)
\end{itemize}

Das Belohnungssystem würde Balance-Erhaltung, Spurgenauigkeit und Energieeffizienz optimieren:
\begin{equation}
R_t = w_1 \cdot R_{balance} + w_2 \cdot R_{tracking} + w_3 \cdot R_{energy} + w_4 \cdot R_{safety}
\end{equation}

\subsubsection{Hierarchisches Reinforcement Learning}
Ein hierarchischer Ansatz könnte verschiedene Abstraktionsebenen implementieren:
\begin{itemize}
    \item \textbf{High-Level Policy}: Pfadplanung und Geschwindigkeitsstrategie
    \item \textbf{Mid-Level Policy}: Kurvenstrategie und Hindernisumfahrung  
    \item \textbf{Low-Level Policy}: Direkte Motor-Ansteuerung für Balance
\end{itemize}

\subsubsection{Transfer Learning und Sim-to-Real}
Die Übertragung der in Webots trainierten RL-Modelle auf reale Hardware erfordert Domain Adaptation-Techniken. Ansätze wie Domain Randomization während des Trainings und progressive Realitätsanpassung könnten die Sim-to-Real-Lücke überbrücken.

\subsection{Hardware-in-the-Loop für realitätsnahe Validierung}

Die Validierung des entwickelten Systems erfordert eine schrittweise Annäherung an reale Bedingungen durch Hardware-in-the-Loop (HIL) Simulationen.

\subsubsection{Embedded Controller Integration}
Die Portierung der C-basierten Balance-Regelung auf Embedded-Systeme wie BeagleBone Black oder Raspberry Pi würde realistische Rechenzeit- und Ressourcenbeschränkungen einführen. Kritische Aspekte umfassen:

\begin{itemize}
    \item \textbf{Real-Time Constraints}: Garantierte Ausführungszeiten für die 500Hz Balance-Regelschleife
    \item \textbf{Interrupt-Handling}: Prioritätsbasierte Behandlung von Sensor-Interrupts
    \item \textbf{Memory Management}: Optimierte Datenstrukturen für begrenzte RAM-Ressourcen
    \item \textbf{Power Management}: Energieeffiziente Algorithmen für Batteriebetrieb
\end{itemize}

\subsubsection{Sensorik-Hardware-Integration}
Die Integration realer Sensoren würde Rauschen, Latenz und Ausfälle einführen:
\begin{itemize}
    \item \textbf{IMU-Kalibrierung}: Bias-Kompensation und Temperaturkorrekturen
    \item \textbf{Kamera-Latenz}: Bildverarbeitungs-Pipeline mit ~50ms Verzögerung
    \item \textbf{Sensor-Ausfälle}: Graceful Degradation bei Sensor-Fehlern
\end{itemize}

\subsubsection{Aktorik-Hardware-Modellierung}
Realistische Motormodelle würden Trägheit, Backlash und Sättigungseffekte berücksichtigen:
\begin{equation}
\tau_{motor} = K_t \cdot i_{motor} - B \cdot \omega_{motor} - \tau_{friction}(\omega_{motor})
\end{equation}

\subsection{Systemintegration und Validierung}

\subsubsection{Modular Testing Framework}
Ein systematisches Test-Framework würde verschiedene Komplexitätsstufen ermöglichen:
\begin{enumerate}
    \item \textbf{Unit Tests}: Einzelne Algorithmus-Komponenten
    \item \textbf{Integration Tests}: Sensor-Fusion und Controller-Interaktion  
    \item \textbf{System Tests}: Vollständige HIL-Simulation
    \item \textbf{Field Tests}: Reale Hardware-Validierung
\end{enumerate}

\subsubsection{Performance Metriken}
Quantitative Bewertungskriterien würden umfassen:
\begin{itemize}
    \item \textbf{Balance-Performance}: RMS Roll-Winkel, Zeit bis zum Umfallen
    \item \textbf{Tracking-Genauigkeit}: Laterale Abweichung, Pfadverfolgungsfehler
    \item \textbf{Energieeffizienz}: Verbrauch pro Kilometer, Regenerationsrate
    \item \textbf{Robustheit}: Performance unter Störungen und Sensorausfällen
\end{itemize}

\subsection{Fazit und Ausblick}

Die vorgeschlagenen Erweiterungen würden das selbstbalancierende Fahrrad von einem Simulationsprototyp zu einem robusten, adaptiven System entwickeln. Die Kombination aus fortgeschrittener Sensor-Fusion, lernfähigen Regelungsalgorithmen und realitätsnaher Hardware-Integration stellt eine natürliche Evolution des aktuellen Ansatzes dar. 

Besonders vielversprechend ist die Möglichkeit, durch Reinforcement Learning nicht nur die Regelungsperformance zu optimieren, sondern auch neue Fahrstrategien zu entdecken, die über klassische Regelungsansätze hinausgehen. Die systematische HIL-Validierung würde dabei sicherstellen, dass die entwickelten Algorithmen auch unter realen Bedingungen zuverlässig funktionieren.

Diese Forschungsrichtungen adressieren sowohl die technischen Herausforderungen der Systemrealisierung als auch die wissenschaftlichen Fragen der optimalen Regelung komplexer, unteraktuierter Systeme und tragen somit zur Weiterentwicklung autonomer Fahrzeugsysteme bei.
